\section{Metric space topology: open, closed, continuity, density, equivalence (Section 4)}

\begin{plain}
``Topology'' is everything that can be described using \emph{open sets} alone: closed sets,
interior, closure, boundary, density, continuity, convergence. Two metrics that have the same
open sets (equivalent metrics) have the same topology, even if their distances are very different.
This is the most heavily examined Section. Every past paper had a full question on it.
\end{plain}

\topic{Open sets}
\begin{defn}{open set}
$U\subset E$ is \textbf{open} if $\forall p\in U\ \exists r>0$ with $B_r(p)\subset U$.
\emph{``Wherever you stand in $U$, there is a ball around you that stays inside $U$.''}
\end{defn}

\begin{exam}{open balls are open (2025 Q3(b) style)}
Prove that $B_R(a)$ is an open set.
\end{exam}
\begin{answer}
Let $p\in B_R(a)$ and set $r=R-d(p,a)>0$ (this is the ``spare radius''). If $x\in B_r(p)$, then
by (MS3) $d(x,a)\le d(x,p)+d(p,a)<r+d(p,a)=R$, so $x\in B_R(a)$. Hence $B_r(p)\subset B_R(a)$,
and $B_R(a)$ is open.
\end{answer}
\begin{center}
\begin{tikzpicture}[scale=1.1]
\draw[thick,navy] (0,0) circle (2); \fill (0,0) circle (1.2pt) node[below left]{$a$};
\fill (1.1,0.5) circle (1.2pt) node[below]{$p$};
\draw[thick,teal,dashed] (1.1,0.5) circle ({2-sqrt(1.1^2+0.5^2)});
\draw[->] (0,0)--(-1.41,1.41) node[midway,left]{$R$};
\draw[->,teal] (1.1,0.5)--++(0,{2-sqrt(1.1^2+0.5^2)}) node[right,pos=.7]{\small $r=R-d(p,a)$};
\end{tikzpicture}
\end{center}

\begin{thm}{the family of open sets $\mathcal O(E)$}
(a) $\emptyset$ and $E$ are open. (b) Any union of open sets is open.
(c) A \emph{finite} intersection of open sets is open (take $t=\min\{r_1,\dots,r_n\}$).
Every open set is a union of open balls: $U=\bigcup_{x\in U}B_{r(x)}(x)$.
\end{thm}
\textbf{Infinite intersections can fail}: $\bigcap_{n}\oi{-\frac1n}{\frac1n}=\{0\}$ is not open
in $\R$.

\begin{exam}{Exercises 4.1 and 4.2}
(a) What are the open sets in a discrete metric space? (b) If $V$ is open and $F$ finite, show
that $V\setminus F$ is open.
\end{exam}
\begin{answer}
(a) \emph{Every} subset $U$ is open: for $p\in U$, $B_1(p)=\{p\}\subset U$.\\
(b) Let $p\in V\setminus F$. Choose $r>0$ with $B_r(p)\subset V$. If $F=\emptyset$ we are done.
Otherwise $F=\{a_1,\dots,a_m\}$, and $\delta=\min_i d(p,a_i)>0$ because $p\notin F$ and the
minimum is over finitely many positive numbers. Let $t=\min\{r,\delta\}$. If $x\in B_t(p)$ then
$x\in V$, and $x\ne a_i$ because $d(x,p)<\delta\le d(a_i,p)$. So $B_t(p)\subset V\setminus F$.
(``Finite'' matters: $\R\setminus\{1/n:n\in\N\}$ is not open, since $0$ has no such ball.)
\end{answer}

\topic{Interior}
\begin{defn}{interior}
$p$ is an \textbf{interior point} of $A$ if $\exists r>0$ with $B_r(p)\subset A$.
$\inte A$ is the set of interior points. $\inte A=\bigcup\{U\subset A: U\text{ open}\}$, the
\emph{largest open subset} of $A$. So $\inte A$ is open, and $A$ is open $\iff A=\inte A$.
\end{defn}
Example: $\inte\Q=\emptyset$ in $\R$, since every interval $\oi{x-r}{x+r}$ contains an
irrational number.

\topic{Continuity}
\begin{defn}{continuity in metric spaces}
$f:E\to E'$ is \textbf{continuous at $p$} if $\forall\varepsilon>0\ \exists\delta>0$ with
$f(B_\delta(p))\subset B_\varepsilon(f(p))$. It is \textbf{sequentially continuous at $p$} if
$x_n\to p\Rightarrow f(x_n)\to f(p)$. These are equivalent. $f$ is continuous if it is continuous
at every point.
\end{defn}

\begin{thm}{global continuity via open sets (very frequently examined)}
$f:E\to E'$ is continuous $\iff$ $f^{-1}(V)$ is open in $E$ for every open $V\subset E'$
$\iff$ $f^{-1}(C)$ is closed in $E$ for every closed $C\subset E'$.
\end{thm}
\begin{answer}
\textbf{($\Rightarrow$)} Let $V\subset E'$ be open and $p\in f^{-1}(V)$. Then $f(p)\in V$, so
there is $\varepsilon>0$ with $B_\varepsilon(f(p))\subset V$. By continuity there is $\delta>0$
with $f(B_\delta(p))\subset B_\varepsilon(f(p))\subset V$. Hence $B_\delta(p)\subset f^{-1}(V)$,
so $f^{-1}(V)$ is open.\\
\textbf{($\Leftarrow$)} Let $p\in E$, $\varepsilon>0$. The ball $V=B_\varepsilon(f(p))$ is open,
so $f^{-1}(V)$ is open and contains $p$. Choose $\delta>0$ with $B_\delta(p)\subset f^{-1}(V)$.
Then $f(B_\delta(p))\subset V$, which is continuity at $p$.\\
\textbf{Closed version (Exercise 4.5).} $f^{-1}(E'\setminus C)=E\setminus f^{-1}(C)$, and a set
is closed iff its complement is open.
\end{answer}

\begin{tip}
Continuity is about \emph{preimages}. The \emph{image} of an open set under a continuous map
need not be open: the constant map $f\equiv0$ on $\R$ sends $\oi{-1}1$ to $\{0\}$, which is not
open. Likewise the image of a closed set need not be closed ($t\mapsto\frac1{1+t}$ sends
$\lo0\infty$ to $\ro01$).
\end{tip}

\textbf{Lipschitz $\Rightarrow$ continuous}: $f(B_r(p))\subset B_{Lr}(f(p))$ with $L=\Lip f$;
take $\delta=\varepsilon/L$. Compositions of continuous maps are continuous (Exercise 6.6:
$(g\circ f)^{-1}(W)=f^{-1}(g^{-1}(W))$).

\topic{Closed sets, closure}
\begin{thm}{point of closure: three equivalent conditions}
For $A\subset E$, $p\in E$, the following are equivalent:
(i) $\forall r>0$, $B_r(p)\cap A\ne\emptyset$; (ii) there is a sequence $(a_n)$ in $A$ with
$a_n\to p$; (iii) $\dist_A(p)=0$. Such a $p$ is a \textbf{point of closure} of $A$, and $\cl A$
is the set of all of them.
\end{thm}
\begin{defn}{closed set}
$A$ is \textbf{closed} if $\cl A=A$, i.e.\ \emph{whenever a sequence in $A$ converges in $E$,
its limit lies in $A$.}
\end{defn}

\begin{thm}{closed $\iff$ complement open}
$A$ is closed $\iff$ $E\setminus A$ is open.
\end{thm}
\begin{answer}
\textbf{Model proof (sequence version, 2025 Q3(e) style).}
($\Leftarrow$) Suppose $E\setminus A$ is open. Let $a_n\in A$ with $a_n\to p$. If $p\notin A$,
there is $r>0$ with $B_r(p)\subset E\setminus A$. Since $a_n\to p$, eventually $a_n\in B_r(p)$,
so $a_n\notin A$, a contradiction. So $p\in A$ and $A$ is closed.\\
($\Rightarrow$) Suppose $A$ is closed and $E\setminus A$ is not open. Then some
$p\in E\setminus A$ has $B_{1/n}(p)\cap A\ne\emptyset$ for every $n$. Pick $a_n$ in this
intersection. Then $a_n\to p$, $a_n\in A$, so $p\in A$ by closedness, a contradiction.
\end{answer}

\begin{thm}{the family of closed sets}
$\emptyset$, $E$ are closed; arbitrary \emph{intersections} and \emph{finite} unions of closed
sets are closed. $\cl A$ is closed and $\cl A=\bigcap\{F\supset A: F\text{ closed}\}$, the
\emph{smallest closed superset} of $A$. Also $\cl{\cl A}=\cl A$.
\end{thm}
\textbf{Infinite unions can fail}: $\bigcup_n[\frac1n,1]=\ro01$ is not closed.

\textbf{Complement identities}: $(\cl A)\cpl=\inte(A\cpl)$ and $(\inte A)\cpl=\cl{A\cpl}$.

\begin{thm}{closures of balls}
$\cl{B_r(a)}\subset B_r[a]$. Equality holds in normed spaces but can fail in general: in a
discrete space with $\ge2$ points, $\cl{B_1(a)}=\{a\}$ but $B_1[a]=E$.
\end{thm}
\begin{answer}
If $p\in\cl{B_r(a)}$, pick $x_n\in B_r(a)$ with $x_n\to p$. Then $d(x_n,a)\to d(p,a)$, and since
$d(x_n,a)<r$ for all $n$, order preservation gives $d(p,a)\le r$. \emph{(Only $\le$: strict
inequality is lost in the limit.)} Closed balls are closed by the same argument, and spheres are
closed too.
\end{answer}

\begin{exam}{closure of unions and intersections (2023 Q2(b) style)}
Prove $\cl{A\cup B}=\cl A\cup\cl B$. Give an example with $\cl{A\cap B}\subsetneq\cl A\cap\cl B$.
\end{exam}
\begin{answer}
$A\subset A\cup B$ gives $\cl A\subset\cl{A\cup B}$; similarly for $B$; so
$\cl A\cup\cl B\subset\cl{A\cup B}$. Conversely $\cl A\cup\cl B$ is closed (a finite union of
closed sets) and contains $A\cup B$, so it contains the smallest closed superset
$\cl{A\cup B}$. For the example, take $A=\oi01$, $B=\oi12$ in $\R$: $\cl{A\cap B}=\cl\emptyset
=\emptyset$, but $\cl A\cap\cl B=[0,1]\cap[1,2]=\{1\}$.
\end{answer}

\topic{Boundary}
\begin{defn}{boundary}
$\partial A=\cl A\setminus\inte A=\cl A\cap\cl{A\cpl}$. It is closed, and
$E=\inte A\ \sqcup\ \partial A\ \sqcup\ \inte(A\cpl)$ (disjoint union).
\end{defn}

\begin{steps}{finding interior, closure and boundary}
\begin{enumerate}
\item Draw the set. \textbf{Interior}: points with a whole ball inside $A$.
\item \textbf{Closure}: add all limits of sequences in $A$.
\item \textbf{Boundary}: closure minus interior. Check each claim with a ball or a sequence.
\end{enumerate}
\end{steps}

\begin{center}
\renewcommand{\arraystretch}{1.25}
\begin{tabular}{@{}lllll@{}}\toprule
Set $A$ (in $\R$ unless stated) & $\inte A$ & $\cl A$ & $\partial A$ & open/closed?\\\midrule
$\lo01$ & $\oi01$ & $[0,1]$ & $\{0,1\}$ & neither\\
$\lo01\cup\{2\}$ & $\oi01$ & $[0,1]\cup\{2\}$ & $\{0,1,2\}$ & neither\\
$\Q$ & $\emptyset$ & $\R$ & $\R$ & neither\\
$\{1/n:n\in\N\}$ & $\emptyset$ & $A\cup\{0\}$ & $A\cup\{0\}$ & neither\\
$\Z$ & $\emptyset$ & $\Z$ & $\Z$ & closed\\
$\R$, $\emptyset$ & itself & itself & $\emptyset$ & both\\
any $A$ in a discrete space & $A$ & $A$ & $\emptyset$ & both\\
$\{(x,y):x^2+y^2<1\}\subset\R^2$ & itself & closed disc & unit circle & open\\
\bottomrule\end{tabular}
\end{center}

\topic{Density and separability}
\begin{defn}{dense, separable}
$D$ is \textbf{dense} in $E$ if $\cl D=E$, i.e.\ every ball $B_R(p)$ meets $D$, i.e.\ every
$p\in E$ is a limit of a sequence in $D$. $E$ is \textbf{separable} if it has a
\emph{countable} dense subset.
\end{defn}
Examples: $\Q$ and $\R\setminus\Q$ are both dense in $\R$ (and disjoint). $\Q^N$ is dense in
$\ell^p(N)$, so $\R^N$ is separable. $C(I;\R)$ with $\dsup$ is separable (polynomials with
rational coefficients; or Exercise 4.11).

\begin{exam}{dense intersect open dense (2023 Q3(a) / 2024 Q2(d))}
Let $U$ be open and dense and $A$ dense in $E$. Prove that $U\cap A$ is dense. Show this can fail
if $U$ is not open.
\end{exam}
\begin{answer}
Let $p\in E$, $R>0$. We must show $B_R(p)\cap U\cap A\ne\emptyset$.
\begin{enumerate}
\item $U$ is dense, so there is $u\in B_R(p)\cap U$.
\item $B_R(p)\cap U$ is open (a finite intersection of open sets), so there is $r>0$ with
  $B_r(u)\subset B_R(p)\cap U$.
\item $A$ is dense, so there is $a\in B_r(u)\cap A$.
\item Then $a\in B_R(p)\cap U\cap A$.
\end{enumerate}
\textbf{Failure without openness}: $U=\Q$, $A=\R\setminus\Q$ are both dense in $\R$, but
$U\cap A=\emptyset$.
\end{answer}

\begin{exam}{a non-separable space}
Show that an uncountable set $S$ with the discrete metric is not separable.
\end{exam}
\begin{answer}
Let $D$ be dense. For each $s\in S$, the ball $B_1(s)=\{s\}$ must meet $D$, so $s\in D$. Thus
$D=S$, which is uncountable. So no countable dense subset exists.
\emph{(The same idea shows that $\ell^\infty(\N;\R)$, the bounded sequences with $\dsup$, is not
separable: the uncountably many $0$--$1$ sequences are pairwise at distance $1$, and the balls
$B_{1/2}$ around them are disjoint, so each must contain a different point of $D$.)}
\end{answer}

\begin{exam}{Exercise 4.11: a countable dense set in $C(I;\R)$ (2024 Q2(d)(v))}
Let $A_n$ be the continuous $g$ that are linear on each $[k/2^n,(k+1)/2^n]$ with
$g(j/2^n)\in\Q$, and $A=\bigcup_n A_n$. Show $A$ is countable and dense in $(C(I;\R),\dsup)$.
\end{exam}
\begin{answer}
\textbf{Countable.} $g\in A_n$ is determined by its values at the $2^n+1$ nodes, so
$g\mapsto(g(0),g(2^{-n}),\dots,g(1))$ is injective into $\Q^{2^n+1}$, which is countable. A
countable union of countable sets is countable.\\
\textbf{Dense.} Let $f\in C(I;\R)$, $\varepsilon>0$. By the Uniform Continuity Theorem choose
$n$ with $|f(s)-f(t)|<\varepsilon/3$ when $|s-t|\le2^{-n}$. Choose rationals $q_j$ with
$|q_j-f(j/2^n)|<\varepsilon/3$, and let $g\in A_n$ take these node values. For $t$ in a mesh
interval with nodes $s_0,s_1$, $g(t)$ is a convex combination of $q_0,q_1$, each within
$\varepsilon/3+\varepsilon/3$ of $f(t)$ (since $|q_i-f(t)|\le|q_i-f(s_i)|+|f(s_i)-f(t)|$).
So $|g(t)-f(t)|<2\varepsilon/3$ for all $t$, and $\dsup(f,g)\le2\varepsilon/3<\varepsilon$.
\end{answer}

\begin{thm}{open subsets of $\R$}
Every open $U\subset\R$ is a countable union of open intervals (those with rational endpoints
contained in $U$).
\end{thm}

\topic{Homeomorphisms}
\begin{defn}{homeomorphism}
$f:E\to E'$ is a \textbf{homeomorphism} if it is continuous, bijective, and $f^{-1}$ is
continuous. Equivalently: a continuous bijection that maps open sets to open sets.
\end{defn}
\textbf{Examples}: $\tan:\oi{-\pi/2}{\pi/2}\to\R$ (a bounded space homeomorphic to an unbounded
one). Bi-Lipschitz maps.\\
\textbf{Non-example}: $\id:(\R,d_{disc})\to(\R,|\cdot|)$ is a continuous bijection (every map
\emph{from} a discrete space is continuous, since every preimage is open), but its inverse is not
continuous: $\{0\}$ is open in the discrete space but not in $\R$.

\topic{Equivalent metrics}
\begin{defn}{topological and Lipschitz equivalence}
Metrics $d_1,d_2$ on $S$ are \textbf{topologically equivalent} if
$\id:(S,d_1)\to(S,d_2)$ is a homeomorphism; \textbf{Lipschitz equivalent} if it is
bi-Lipschitz, i.e.\ $\exists m,M>0$ with $m\,d_1\le d_2\le M\,d_1$.
Lipschitz equivalent $\Rightarrow$ topologically equivalent.
\end{defn}
\begin{thm}{tests for topological equivalence}
The following are equivalent: (i) $d_1,d_2$ topologically equivalent; (ii) they have the same
convergent sequences (with the same limits); (iii) they have the same closed sets; (iv) they have
the same open sets; (v) every $d_1$-ball about $p$ contains a $d_2$-ball about $p$, and vice versa.
\end{thm}
\begin{center}
\renewcommand{\arraystretch}{1.2}
\begin{tabular}{@{}lcc@{}}\toprule
Property & top.\ equiv.\ preserves? & Lipschitz equiv.\ preserves?\\\midrule
open and closed sets, convergence, density & yes & yes\\
sequential compactness & yes & yes\\
boundedness & \textbf{no} & yes\\
Cauchy sequences, completeness & \textbf{no} & yes\\\bottomrule
\end{tabular}
\end{center}

\begin{thm}{$\alpha\circ d$ is topologically equivalent to $d$}
If $\alpha$ is definite, increasing, subadditive and \emph{continuous at $0$}, then $\alpha\circ
d$ is a topologically equivalent metric. If moreover $\alpha$ is strictly increasing near $0$,
then $d$ and $\alpha\circ d$ have the same Cauchy sequences. Examples: $cs$, $\min\{s,1\}$,
$s/(1+s)$, $\tan^{-1}s$, $s^{1/p}$. \textbf{Corollary}: every metric is topologically equivalent
to a bounded metric with the same Cauchy sequences.
\end{thm}

\begin{exam}{truncated and arctangent metrics (2024 Q3(b)(c) style)}
(a) Show that $d'=\min\{d,1\}$ has the same convergent and the same Cauchy sequences as $d$.
(b) Show $d'$ and $d$ need not be Lipschitz equivalent.
(c) Show $d_{NS}(x,y)=|\tan^{-1}x-\tan^{-1}y|$ is topologically, but not Lipschitz, equivalent
to the usual metric on $\R$, and that $(\R,d_{NS})$ is not complete.
\end{exam}
\begin{answer}
(a) $d'\le d$, so $d$-convergent/Cauchy $\Rightarrow$ $d'$-convergent/Cauchy. Conversely, given
$0<\varepsilon<1$: once $d'(x_n,x_m)<\varepsilon<1$, the minimum must equal $d(x_n,x_m)$, so
$d(x_n,x_m)<\varepsilon$. The same holds for $d(x_n,a)$.\\
(b) On $\R$ with $d=|\cdot|$: if $d\le M\,d'$, then taking $x=0$, $y=n$ gives $n\le M$ for all
$n$, which is impossible. (In general: Lipschitz equivalence preserves boundedness, $d'$ is
bounded, and $d$ is not.)\\
(c) $\tan^{-1}:\R\to\oi{-\pi/2}{\pi/2}$ is a homeomorphism, so $x_n\to x$ in $|\cdot|$ $\iff$
$\tan^{-1}x_n\to\tan^{-1}x$ $\iff$ $x_n\to x$ in $d_{NS}$. Not Lipschitz equivalent: $d_{NS}<\pi$
is bounded, the usual metric is not. Not complete: $x_n=n$ is $d_{NS}$-Cauchy because
$\tan^{-1}n\to\pi/2$, so $d_{NS}(n,m)\le|\tan^{-1}n-\frac\pi2|+|\frac\pi2-\tan^{-1}m|\to0$. If
$n\to a$ in $d_{NS}$ then $\tan^{-1}n\to\tan^{-1}a$, so $\tan^{-1}a=\pi/2$, which is impossible.
\textbf{Moral: completeness is not a topological property.}
\end{answer}

\begin{exam}{Exercise 4.8: $d_{TGV}$ on $\R$}
$d_{TGV}(x,y)=|x|+|y|$ for $x\ne y$. Let $f=\id:(\R,|\cdot|)\to(\R,d_{TGV})$ and $g=f^{-1}$.
Show (a) $g$ is Lipschitz; (b) $f$ is continuous at $0$; (c) $f$ is not continuous at $a\ne0$.
\end{exam}
\begin{answer}
(a) $|x-y|\le|x|+|y|=d_{TGV}(x,y)$ for $x\ne y$, so $\Lip g\le1$.\\
(b) $d_{TGV}(x,0)=|x|$, so $|x-0|<\varepsilon\Rightarrow d_{TGV}(x,0)<\varepsilon$ ($\delta=\varepsilon$).\\
(c) Take $x_n=a+\frac1n\to a$ in $|\cdot|$. Then $d_{TGV}(x_n,a)=|x_n|+|a|\to2|a|\ne0$, so
$f(x_n)\not\to f(a)$. Hence $d_{TGV}$ is not topologically equivalent to the usual metric.
(Indeed $\{a\}=B^{TGV}_{|a|}(a)$ is $d_{TGV}$-open.)
\end{answer}

\textbf{All $\ell^p$ metrics on $\R^N$ are Lipschitz equivalent}, since
$\norm x_\infty\le\norm x_p\le N^{1/p}\norm x_\infty$. So in $\R^N$, open, closed, bounded,
convergent, Cauchy and compact do not depend on $p$, and convergence is coordinatewise.

\topic{Metric subspaces: ``open in $A$'' versus ``open in $E$''}
For $A\subset E$ with the restricted metric: $B^A_r(a)=A\cap B_r(a)$.
\begin{thm}{relatively open and closed sets}
For $X\subset A$: $X$ is open in $A$ $\iff$ $X=A\cap U$ for some $U$ open in $E$;
$X$ is closed in $A$ $\iff$ $X=A\cap F$ for some $F$ closed in $E$.
\end{thm}
\begin{answer}
\textbf{Proof of the open version (Exercise 4.10; 2022 Q2(d)).}
($\Leftarrow$) Let $v\in X=A\cap U$. Choose $r$ with $B_r(v)\subset U$. Then
$B^A_r(v)=A\cap B_r(v)\subset A\cap U=X$.\\
($\Rightarrow$) For each $v\in X$ choose $r_v>0$ with $B^A_{r_v}(v)\subset X$. Let
$U=\bigcup_{v\in X}B_{r_v}(v)$, which is open in $E$ as a union of balls. Each $v\in X$ is in
$A\cap U$. If $x\in A\cap U$ then $x\in A\cap B_{r_v}(v)=B^A_{r_v}(v)\subset X$ for some $v$.
So $X=A\cap U$.
\end{answer}

\begin{exam}{clopen in the subspace, neither in $E$ (2024 Q2(c)(ii))}
Give $A\subset E$ and $R\subset A$ such that $R$ is open and closed in $A$, but neither open nor
closed in $E$.
\end{exam}
\begin{answer}
$E=\R$, $A=R=\lo01$. The whole space is always open and closed in itself. In $\R$: not open,
since every ball about $0$ contains negative numbers; not closed, since $1-\frac1n\in R$
converges to $1\notin R$.
Another: in $A=\Z\subset\R$, $B^{\Z}_1(k)=\{k\}$, so every subset of $\Z$ is open and closed in
$\Z$.
\end{answer}

\begin{exam}{identity theorem (Exercise 6.7)}
Let $f,g:E\to E'$ be continuous. Show $H=\{x: f(x)=g(x)\}$ is closed, and that if $f=g$ on a
dense set then $f=g$.
\end{exam}
\begin{answer}
If $x_n\in H$, $x_n\to x$, then $f(x)=\lim f(x_n)=\lim g(x_n)=g(x)$ by sequential continuity and
uniqueness of limits, so $x\in H$. If $D\subset H$ is dense then $E=\cl D\subset\cl H=H$.
\end{answer}

\begin{tip}
\begin{itemize}
\item Always say \emph{where} a set is open or closed: ``open in $A$'' or ``open in $E$''.
\item To prove ``closed'' you may use sequences \emph{or} the complement. Pick the shorter route.
\item Many proofs need a radius $r$. Draw the picture and use the spare radius
  $R-d(p,a)$.
\item ``Dense'' means meets every ball. Prove it by fixing an arbitrary ball.
\item Completeness and boundedness are \emph{not} preserved by mere topological equivalence.
\end{itemize}
\end{tip}
